\section{Discussion}
A fitted population field brought day-1 cells closer to the observed day-7 distribution in this acute human wound cohort. Improvement was present when day 7 was withheld globally and under a separate protocol that withheld each donor. These comparisons support distributional reconstruction for the measured interval in three healthy volunteers. Their distinction remains essential: transfer to a new donor can use day-7 observations from other people, whereas global timepoint holdout cannot. Neither task evaluates diabetic-foot-ulcer prognosis or the subsequent course of an individual chronic wound.

The mathematical interpretation is narrower than a recovered biological trajectory. Training samples independent cells from the endpoints of a donor's interval and regresses their linear-path velocities. It imposes neither a cell-lineage match nor a minimum-cost coupling. At the ideal regression solution, probability conservation concerns that specified path and sampling measure. It does not guarantee the distribution at an unobserved biological time, and empirical endpoints alone do not supply all the regularity needed for an exact characteristic-flow interpretation. This distinction is consistent with the fundamental ambiguity of dynamics inferred from snapshots and with the stronger assumptions required by identification theory \citep{weinreb2018,lavenant2021}.

The time-axis comparison exposes a consequential modeling choice. Linear days performed poorly at the target despite favorable results under several alternatives. Changing the axis here changes the straight bridges, their velocity scales and the optimization problem; the experiment is not a coordinate transformation of one already fitted field. Compression of the earliest interval is a possible contributor, but its separate effect was not isolated. The comparison therefore supports sensitivity to the chosen clock, not a preferred biological clock. Selecting an axis after inspecting the same held-out target also requires independent evaluation before claiming predictive validation \citep{cawley2010}.

Conditioning on a donor's source distribution improved training-donor fit but did not improve the overall held-out mean relative to the shared field. This separates additional fitting capacity from transfer performance. A shared field can give different velocities at different latent positions, so differences between donor-average displacements do not by themselves prove that donor context is necessary. The non-neural predictors also improved on unchanged source, indicating that part of the available structure can be captured by simple displacement rules. Broader comparisons of context-dependent and transport-based predictors should retain these simple controls \citep{bunne2023}.

The training and evaluation measures deserve explicit attention. Donor--segment cycling approximately balances donors during fitting, whereas the original pooled error weights their cells by yield. Re-evaluation of the same fixed prediction with equal donor mass preserved the improvement, and each donor separately improved under that pooled time-holdout field. The pooled result is therefore not explained solely by the observed cell-count imbalance. This sensitivity calculation does not refit a donor-balanced scaler, measure performance in a new donor, or remove dependence on the transferred expression coordinates.

Numerical integration and evaluation sampling were comparatively stable for the tested fixed field. Increasing RK4 resolution produced negligible changes, and repeated cell subsampling preserved the improvement. These checks reduce concern about those specific computational choices while leaving representation and training uncertainty open. Individual training configurations varied substantially across seeds. The aggregate gain from adding a second training donor persisted, but two training sizes cannot identify a learning-curve plateau. Overlapping training sets also correlate held-out errors, limiting the interpretation of a bootstrap interval based on only three donor blocks \citep{bengio2004}.

The principal limit is biological replication. Four times from three healthy donors describe a restricted acute-wound setting, and many cells do not increase the number of independent people \citep{zimmerman2021}. Each mouse arm has one animal per time point, so differences between times also include differences between animals. Incomplete gene coverage and variable performance further restrict cross-species interpretation. The next biological test should expand independent-donor evaluation under a frozen representation, clock and evaluation rule. An explicitly longitudinal diabetic-foot-ulcer cohort with authoritative outcomes would then be needed to test clinical staging or prognosis.
